\clearpage\begingroup\NoSourceHead\SourceLeaf{530b}
\renewcommand{\thefootnote}{\arabic{footnote})}\setcounter{footnote}{30}
\ArithmeticHead{73}{R}{DER SATZ ÜBER DIE TOTALKRÜMMUNG.}
\begin{center}\textls[100]{2. Der \textsc{Gauss}sche Satz über die Totalkrümmung}\\
\textls[100]{und seine von \textsc{Bonnet} gegebene Verallgemeinerung.}\end{center}
\noindent\hspace*{\parindent}17. Im letzten Abschnitt\footnote{\textsc{Gauss}, loc. cit. 24), S. 394—395.} der in der vorigen Nummer genannten nachgelassenen Arbeit gibt \textsc{Gauss} eine merkwürdige Umformung des von ihm für
die geodätische Krümmung gefundenen Ausdrucks, die den eigentlichen Kern
nicht nur des \textsc{Gauss}schen Satzes über die Totalkrümmung, sondern auch seiner
später von \textsc{Bonnet}\footnote{O. \textsc{Bonnet}, \textit{Mémoire sur la théorie générale des surfaces}, Journal de l'École Polytechnique
Cahier 32 (1848), S. 131.} gegebenen Verallgemeinerung enthält. Der ursprünglich
von \textsc{Gauss} für die geodätische Krümmung abgeleitete Ausdruck hat die folgende Gestalt
\begin{equation*}\tag{44}
\begin{aligned}
K_g\,ds&=\frac{\sqrt{EG-F^2}(p'\,dq'-q'\,dp')}{Ep'^2+2Fp'q'+Gq'^2}\\
&\quad+\frac{\Phi(p,q,p',q')\,dt}{(Ep'^2+2Fp'q'+Gq'^2)\sqrt{EG-F^2}},
\end{aligned}
\end{equation*}
wobei $\Phi$ eine kubische Form der Variabeln $p'$, $q'$ ist, deren Koeffizienten
rationale Funktionen von $E$, $F$, $G$ und ihren ersten Ableitungen nach $p$ und
$q$ sind. Nun macht \textsc{Gauss} die wichtige Bemerkung, dass für das Differential
des Winkels $\theta$ zwischen der positiven $p$-Richtung und der positiven Tangente
der Kurve in dem betrachteten Punkt ein genau analoger Ausdruck gilt, der
sich von dem obigen nur dadurch unterscheidet, dass darin die Funktion $\Phi$
durch eine andere kubische Form von $p'$, $q'$ ersetzt ist, sodass also bei Subtraktion beider Ausdrücke das erste, die zweiten Ableitungen von $p$ und $q$
enthaltende Glied herausfällt und man ein Resultat der folgenden Form erhält
\begin{equation*}\tag{45}
K_g\,ds-d\theta=\frac{\Psi(p,q,p',q')\,dt}{(Ep'^2+2Fp'q'+Gq'^2)\sqrt{EG-F^2}},
\end{equation*}
wo $\Psi$ eine neue kubische Form von $p'$, $q'$ bedeutet.
\par\noindent\hspace*{\parindent}Darüber hinaus konstatiert nun \textsc{Gauss} die überraschende Tatsache, dass
\textls[100]{die kubische Form} $\Psi$ \textls[100]{durch} $Ep'^2+2Fp'q'+Gq'^2$ \textls[100]{teilbar ist}, sodass das
Resultat schliesslich die Form annimmt
\begin{equation*}\tag{46}
K_g\,ds-d\theta=P\,dp+Q\,dq,
\end{equation*}
\par\noindent\hbox to\linewidth{\small X 2 Abh. 5.\hfill 10\hspace*{5em}}\par
\clearpage\SourceLeaf{531}
\ArithmeticHead{74}{L}{BOLZA, GAUSS UND DIE VARIATIONSRECHNUNG.}
\noindent wo $P$, $Q$ die folgenden Funktionen von $p$, $q$ sind
\[
P=\frac{EF_p-\frac12EE_q-\frac12FE_p}{E\sqrt{EG-F^2}},\quad
Q=\frac{\frac12EG_p-\frac12FE_q}{E\sqrt{EG-F^2}}.
\]
\par\noindent\hspace*{\parindent}In dieser Formel (46) ist aber in der Tat der \textsc{Gauss}sche Satz über die
Totalkrümmung, und zwar gleich in der \textsc{Bonnet}schen Verallgemeinerung enthalten. Denn für den speziellen Fall eines Systems von geodätischen Parallel-
oder Polarkoordinaten, für das $ds^2=dp^2+m^2dq^2$ ist, nimmt (46) die einfache
Form an
\[
d\theta-K_g\,ds=-\frac{\partial m}{\partial p}\,dq.
\]
\par\noindent\hspace*{\parindent}Integriert man diese Gleichung in positivem Sinn um die Begrenzung $\mathfrak C$
eines einfach zusammenhängenden Stückes $\sigma$ der gegebenen Fläche herum, so
erhält man den Satz von der Totalkrümmung in der \textsc{Bonnet}schen Verallgemeinerung
\begin{equation*}\tag{47}
\int K\,d\sigma=\int\limits_{\mathfrak C}d\theta-\int\limits_{\mathfrak C}K_g\,ds,
\end{equation*}
unter $K$ das Krümmungsmass der Fläche verstanden\footnote{Auch ohne Benutzung eines speziellen Koordinatensystems $p$, $q$ erhält man direkt aus (46) durch
den angegebenen Integrationsprozess die Gleichung (47), wenn man von dem bei L. \textsc{Bianchi}, \textit{Vorlesungen über
Differentialgeometrie} (2. Aufl. 1910), S. 67, Gleichung (17) gegebenen Ausdruck für das Krümmungsmass Gebrauch macht.}. Es ist kaum anzunehmen, dass \textsc{Gauss} dieser naheliegende Zusammenhang entgangen sein sollte,
und so darf man wohl annehmen, dass \textsc{Gauss} die \textsc{Bonnet}sche Verallgemeinerung
seines Satzes bereits gekannt hat, wie dies auch schon R. v. \textsc{Lilienthal}\footnote{\textit{Encyklopädie der Mathematischen Wissenschaften}, Bd. III, D 3, S. 134.} und
P. \textsc{Stäckel}\footnote{P. \textsc{Stäckel}, \textit{Gauss als Geometer} (1918), Werke X 2, Abth. IV, Nr. 32, Materialien für eine wiss.
Biogr. von Gauss, Heft V, S. 126.} aus dem Umstand geschlossen haben, dass \textsc{Gauss} den Ausdruck
$K_g\,ds$ als Differential $df$ schreibt, und als »Differential der Seitenkrümmung«
bezeichnet.
\clearpage\endgroup
